\section{Scope}

GaitProtector \cite{duan2026gaitprotector} is a targeted gait de-identification method that relies on adversarial optimization of latent diffusion features 
to perturb gait patterns in silhouette videos. GaitProtector is able to produce natural looking de-identified videos that closely preserve the 
spatial structure of the original videos while impeding identifications of subjects through gait detectors. Such spatial structure is so well preserved that
identifying gait patterns are removed without disturbing signal that can lead to diagnosis of scoliosis and other medical conditions, enabling de-identification
of medical datasets.

While GaitProtector has a satisfactory performance on silhouette videos of single individuals extracted from a constrained dataset, it was not designed to work on RGB videos and was
not tested by its authors on unconstrained data. The goal of this study is to extend the GaitProtector approach to RGB footage, test it on unconstrained data, 
and evaluate its future integration in pipelines able to de-identify multiple subjects at once. This experiment is part of a broader effort to make gait identification
and deidentification techniques more robust and suitable for real-world data \cite{fan2023opengait}.

In particular, we hope to achieve the following incremental goals:
\begin{enumerate}
    \item \textbf{Replicate the original GaitProtector paper.} This step is necessary because the authors of GaitProtector did not share any of the original code.
    \item \textbf{Adapt GaitProtector to work on constrained RGB data.} The original paper uses VideoFusion \cite{luo2023videofusion} as the main diffusion backbone. Video fusion can take
    RGB input, so we assume the only portion of the GaitProtector pipeline requiring alteration will be the optimization algorithm. We may choose to replace the iterative non-neural
    optimization proposed by GaitProtector with a more sophisticated neural non-iterative module to accommodate for the richer RGB signal.
    \item \textbf{Test the newly developed de-identification system on unconstrained data.} E.g. we may test the model on GREW dataset \cite{Zhu_2021_ICCV} on top of replicating
    the same tests described in the original GaitProtector paper. We plan to retain the same black-box threat evaluation protocol used by the GaitProtector authors.
    \item \textbf{Adapt the newly developed system to multi-individual footage.} This is the most ambitious goal of our study. We plan to try two different approaches. The first
    and most natural approach would entail tracking individuals across the multi-individual video, de-identifying each individual's gait by applying the pipeline to a crop of the original video, and then combining the results
    together. The second more difficult but possibly more elegant approach consists in tracking and de-identifying individuals directly in the latent diffusion space. We do not expect to
    be fully successful with any of these two approaches, but we expect them to set an interesting baseline for work to come.
\end{enumerate}